\chapter{Consciência e Sensibilidade}

\lettrine{P}{or vezes abordamos a meditação} demasiado a partir de um ideal de tentar
controlar a mente e livrarmo-nos de estados mentais indesejáveis. Pode tornar-se
uma obsessão. A meditação pode ser só mais uma coisa que temos que fazer; e esta
atitude mundana normalmente afecta o que estamos a fazer.

Olhemos para a meditação não como algo para nos avaliarmos como pessoas, mas
como uma ocasião ou oportunidade de estarmos conscientes e em paz connosco
próprios e com qualquer estado de espírito ou humor em que possamos estar neste
momento. Aprendamos a ser alguém que está em paz assim como é, ao contrário de
alguém que tenta tornar-se algo ou alcançar um estado que gostaria de ter.

Toda essa forma de pensar é baseada na ilusão. Lembro-me que quando comecei a
meditar na Tailândia, todas as minhas tendências ambiciosas e agressivas
começaram a assumir controlo. A maneira como eu vivia a minha vida afectaria a
forma como eu abordaria a meditação. Então comecei a perceber isso. Comecei a
abrir mão das coisas e até mesmo a aceitar essas tendências e a estar atento à
realidade como é. Quanto mais confiamos nisso, mais rapidamente compreenderemos
o Dhamma, ou a saída do sofrimento.

Aperceba-se de como as coisas afectam a sua mente. Se acabou de chegar do
trabalho ou de sua casa, observe o que isso faz à sua mente. Não critique isso
-- não estamos aqui para culpar ou pensar que existe algo de errado com a nossa
profissão se a nossa mente não estiver tranquila e pura e serena quando vimos
aqui. Mas observe as actividades da vida: ter que falar com as pessoas, ter que
atender telefones, ter que escrever ou atravessar Londres na hora de ponta.
Talvez tenhamos que trabalhar com pessoas de quem não gostamos, em situações
difíceis e agravantes. Apenas observe -- não para criticar, mas apenas para
aceitar que essas coisas têm um efeito sobre nós.

Reconheça que esta é a experiência da consciência e da sensibilidade. É isso que
representa nascer-se como ser humano, não é? Nascemos e temos que viver uma vida
inteira como um ser consciente de uma forma muito sensível. Então, o que nos
afecta, o que vem do mundo objectivo, afectar-nos-á. É simplesmente a realidade
como é. Não há nada de errado com isso.

Só que depois, o ignorante ser humano interpreta tudo muito
individualisticamente e, portanto, tem tendência para levar tudo muito a peito.
É como se eu não me devesse deixar afectar por essas coisas que me atingem. Eu
não devia sentir raiva, aversão, ganância, irritação, frustração, inveja, ciúme,
medo ou ansiedade. Eu não devia sentir estas coisas. Se eu fosse uma pessoa
normal e saudável, não teria nenhum destes problemas. Se eu fosse uma pessoa
normal e saudável, não seria nada sensível. Tal como um rinoceronte, com uma
pele dura onde nada jamais consegue penetrar.

Mas reconheçamos que, sendo humanos, temos estas formas corpóreas extremamente
sensíveis. Não há nada de realmente errado connosco. É simplesmente a realidade
como é. A vida é assim. Vivemos numa sociedade em que a realidade é simplesmente
como é. Morando em Londres ou num subúrbio, nas aldeias ou em qualquer lugar,
podemos passar o nosso tempo a reclamar porque não é perfeito, ou porque existem
muitas coisas irritantes, ou não muito simpáticas em muitos aspectos das nossas
vidas. Mas então, ser-se sensível é assim, certo? Sensibilidade significa isso.
Seja o que for, agradável ou desagradável, aprazível ou doloroso, bonito ou feio
-- nós vamos sentir isso.

Na realidade, a saída do sofrimento é através da consciência. Quando se é
verdadeiramente consciente, não existe eu. Não estamos a recolher experiências
de vida a partir do pressuposto que somos uma pessoa. Podemos tentar tornar-nos
insensíveis, como fecharmos os olhos, colocarmos protectores nos ouvidos e assim
tentamos ser totalmente insensíveis, desligando tudo -- e esse é um tipo de
meditação, privação sensorial. E se nos deixarmos ficar assim por algum tempo,
sentir-nos-emos muito calmos, porque nada é exigido de nós naquele momento. Não
há qualquer tipo de impacto difícil, estimulante, excitante ou frustrante.

Então, se formos conscientes, temos uma consciência da pureza da nossa mente, o
que é uma felicidade. Assim, a nossa verdadeira natureza é feliz, serena e pura.
Mas, se ainda tivermos uma visão errada das coisas, pensamos, ``Eu tenho que ter
uma experiência de privação sensorial o tempo inteiro. Não posso viver mais em
Londres. Até a Sociedade Budista\footnote{O local desta palestra foi na sala de
  meditação da Sociedade Budista em Victoria, Central London.} é muito
barulhenta!''

Se a nossa paz e serenidade dependerem das condições terem de ser de certa
forma, então ficamos muito apegados. Tornamo-nos escravos, queremos controlar as
situações e depois ainda ficamos mais incomodados se algo perturba essas
condições e interfere na nossa paz. ``Tenho que encontrar um lugar -- uma gruta.
Tenho que arranjar o meu próprio tanque de privação sensorial e encontrar a
situação ideal. Estabelecer todas as condições onde consiga manter tudo sob
controlo, para que possa simplesmente permanecer na serenidade abençoada da
pureza da mente.'' Mas então vejamos, essa visão é baseada no desejo, certo? É
uma visão egoísta, um desejo de ter essa experiência porque nos lembramos,
gostámos dela e queremo-la de novo.

Uma vez, num retiro, escutei uma pessoa que estava a ter problemas em engolir.
Eu estava lá sentado e ouvia a pessoa a fazer `glup glup'. Não se ouvia muito,
mas quando estamos apegados ao silêncio total, até um `gulp' nos pode incomodar.
Fiquei bastante irritado e tive vontade de expulsá-la da sala de meditação.
Reflectindo sobre isso, percebi que a culpa estava em mim e não na pessoa.

Mas a consciência e a compreensão do Dhamma permitem que nos adaptemos e
aceitemos a vida -- a experiência total da vida -- sem precisarmos de
controlá-la. Com consciência, não precisamos de nos agarrar a ínfimos aspectos
da realidade que gostamos e depois sentirmo-nos muito ameaçados pelas
possibilidades de podermos ficar separados desses aspectos. A meditação correcta
permite-nos ser muito corajosos e adaptáveis, flexíveis com a nossa vida e com
tudo o que isso implique.

Não temos assim tanto controlo, não é? Por mais que queiramos ser capazes de
controlar nossas vidas, reconhecemos que realmente não temos muito controlo.
Algumas coisas fogem do nosso controlo simplesmente. As coisas acontecem e a Mãe
Natureza tem as suas maneiras de nos informar que não vai simplesmente seguir os
nossos desejos. Depois temos modas e revoluções, mudanças nas condições e
problemas populacionais, e aviões, televisões, tecnologia e poluição -- como
podemos controlar e fazer com que não sejamos afectados por nada disto, ou que
sejamos afectados apenas da maneira que gostamos?

Se passarmos as nossas vidas a tentarmos controlar tudo, apenas aumentamos o
sofrimento. Mesmo que arranjássemos uma forma de controlo sobre as coisas que
queremos, ainda assim seríamos como eu com a pessoa a engolir em seco na sala de
meditação, bem como ficarmos muito irritados quando o vizinho põe o rádio muito
alto ou quando o avião passa muito baixo ou o carro dos bombeiros faz muito
barulho.

Mas uma coisa que temos de reconhecer é: quando temos um corpo, temos que
conviver com ele durante a vida inteira. E estes corpos são formas conscientes e
sensíveis. Isto simplesmente é a realidade como é, o que significa nascer.
Corpos: eles crescem, depois começam a envelhecer. E então temos a velhice e
logo as enfermidades, a doença -- isto faz parte da nossa experiência humana. E
depois a morte. Temos que aceitar a morte e a separação de entes queridos. Isso
acontece com todos nós. A maioria de nós verá os nossos pais falecerem, ou mesmo
os nossos filhos, cônjuge, amigos ou entes queridos. Parte de toda experiência
humana é a experiência de se estar separado de quem gostamos e amamos.

Ao conhecer-se a realidade como é, tornamo-nos perfeitamente capazes de aceitar
a vida e não ficarmos deprimidos e confusos como a vida acaba por ser. Uma vez
entendendo e olhando a vida da forma certa, não criaremos nenhuma visão errada
acerca disso. Não vamos acrescentar medos, desejos, amarguras, ressentimentos e
culpas. Temos a capacidade de aceitar a maneira como a vida nos acontece como
seres individuais. Embora sejamos terrivelmente sensíveis, somos igualmente
sobreviventes resistentes neste universo.

Repare-se onde os seres humanos conseguem viver, como os esquimós no Árctico e
as pessoas nos desertos. Normalmente, existem habitações humanas nos lugares
mais inóspitos deste planeta: quando forçados, conseguimos sobreviver em
qualquer lugar.

Compreender o Dhamma também nos permite ter uma atitude destemida. Começamos a
perceber que podemos aceitar o que quer que aconteça. Não há realmente nada a
temer. E então conseguimos largar a vida, conseguimos segui-la porque já não
esperamos propriamente nada dela e não estamos a tentar controlá-la. Temos a
sabedoria, a consciência e a capacidade de fluir com a corrente, em vez de nos
afogarmos pelas marés e ondas da vida.

Podemos aprender a ficar em silêncio e a ouvir-nos a nós próprios. Usemos
simplesmente a respiração e o corpo. Simplesmente o ritmo natural e a forma como
sentimos o corpo agora. Prestemos atenção ao corpo, porque o corpo é uma
condição na natureza. Nós na realidade não somos o corpo. Já não é a ``minha''
respiração, não é a minha pessoa; nós respiramos mesmo se estivermos dementes ou
doentes -- e se estivermos a dormir, o corpo continua a respirar. O corpo
respira. Do nascimento à morte, estará a respirar. Portanto, a respiração é algo
que usamos como um objecto onde nos focamos, para onde nos viramos. Se pensarmos
muito, os nossos pensamentos tornam-se muito confusos e complicados; mas se
trouxermos a atenção simplesmente para a respiração normal do corpo agora, nesse
momento não estaremos a pensar verdadeiramente -- estamos atentos a um ritmo
natural.

Logo a seguir podemos começar a criar problemas com isso: ``Oh, eu não consigo
concentrar-me na minha respiração, blá, blá, blá, \ldots{}'' Então torna-se um
``eu'' novamente, a tentar estar atento à sua respiração. Mas na realidade, em
qualquer momento em que estamos apenas com a respiração, não há eu. O eu surgirá
quando começarmos a pensar. Quando não estamos a pensar, não existe eu. E quando
estamos conscientes, o pensamento então não virá da perspectiva errada de: ``Eu
sou um eu''. O pensamento então pode ser uma forma de reflexão, uma forma de
focar a atenção no Dhamma, em vez de criar problemas, criticismo e ansiedade
sobre mim e a humanidade.

Contemplemos simplesmente: quando ficamos com raiva temos que pensar, isto está
certo? Se pararmos de pensar, a raiva desaparecerá. Para estarmos irritados,
temos que pensar: ``Ele disse-me isto, como é que ele se atreve. Aquele fulano
nojento!'' Mas se pararmos de pensar e usarmos a respiração, eventualmente a
sensação do corpo que vem com raiva esmorecerá e então não há raiva. Portanto,
se nos sentirmos com raiva, reflictamos simplesmente como é essa sensação
física. É o mesmo com qualquer disposição de ânimo: contemplemos e reflictamos
sobre a disposição de ânimo em que nos encontramos. Trabalhemos simplesmente com
isso -- não para a analisarmos ou criticarmos, mas tão só para reflectirmos
sobre como se encontra a nossa disposição de ânimo.

Às vezes, as pessoas dizem: ``Fico muito confuso quando medito. Como é que me
posso ver livre da confusão?'' Querer ver-se livre da confusão é o problema.
Ficar confuso e não querer, gera mais confusão ainda.

Então, como é a sensação de confusão? Algumas das paixões mais estimulantes que
podemos ter são bastante óbvias. O que temos tendência para não prestar atenção
nenhuma, ou ignorar, são os estados mais subtis como a ligeira confusão,
hesitação ou dúvida, insegurança e ansiedade. E, é claro, um lado de nós só se
quer ver livre disso, lançar isso para longe. ``Como é que me livro disto? Se eu
meditar, como posso ver-me livre dos meus medos e ansiedade?''

Com o entendimento correcto, vemos que o próprio desejo de ``livrarmo-nos de'' é
sofrimento. Conseguimos aguentar o sentimento de insegurança, se soubermos o que
é, o que muda e o que é impermanente. Assim, cada vez mais começamos a
sentir-nos confiantes ao estarmos simplesmente cientes e conscientes, em vez de
tentarmos desenvolver uma prática para nos tornarmos uma pessoa iluminada.
Assume-se que agora não somos iluminados, que temos muitos problemas, que
precisamos de mudar a nossa vida, precisamos de nos tornar diferentes. Não somos
suficientemente bons como estamos agora. Então temos que meditar e, com sorte,
nalgum momento futuro, tornar-nos-emos algo que gostaríamos de ser.

Se nunca discernirmos a ilusão desta forma de pensar, então ela simplesmente
continuará. Nunca nos tornaremos no que realmente deveríamos ser. Não importa
quanto esforço coloquemos nas nossas meditações, depois de anos a tentarmos
iluminar-nos, sempre nos sentiremos um fracasso porque ainda temos a atitude
errada sobre tudo isso.

\photoPage{image-008-forest-path.jpg}{Hammer Wood, Vihara de Cittaviveka, Chithurst}{Foto por Derek Gardiner}
